\documentclass[10pt,a4paper]{amsart}
\usepackage[margin=19mm]{geometry}
\usepackage{amsmath,amssymb,mathtools}
\usepackage[scr=rsfs,cal=euler]{mathalfa}
\usepackage{xfrac}
\usepackage[unicode,hidelinks,bookmarks=false]{hyperref}
\newcommand{\ie}{i.e.\ }
\newcommand{\cf}{cf.\ }
\newcommand{\resp}{respectively\ }
\newcommand{\Subsection}{\subsection}
\providecommand{\nobreakdash}{\nobreak\hbox{-}}
\setlength{\emergencystretch}{2em}
\multlinegap0pt
\usepackage{lmodern}
\usepackage{mathtools}
\hypersetup{
 pdftitle={A Logarithm-Free Critical Endpoint Estimate for the Two-Dimensional Hermite Operator},
 pdfauthor={}
}
\allowdisplaybreaks
\setlength{\parindent}{1.5em}
\setlength{\parskip}{0pt}
\newtheorem{theorem}{Theorem}[section]
\newtheorem{proposition}[theorem]{Proposition}
\newtheorem{lemma}[theorem]{Lemma}
\numberwithin{equation}{section}
\newcommand{\R}{\mathbb R}
\newcommand{\T}{\mathbb T}
\newcommand{\N}{\mathbb N}
\newcommand{\calH}{\mathcal H}
\newcommand{\one}{\mathbf 1}
\newcommand{\dd}{\,\mathrm d}
\newcommand{\norm}[1]{\lVert #1\rVert}
\begin{document}
\title[Sharp Endpoint Eigenfunction Estimates for the Two-Dimension]{Sharp Endpoint Eigenfunction Estimates for the Two-Dimensional Hermite Operator}
\author{Guiyu Xie, Cheng Zhang}
\date{}
\maketitle
\noindent\emph{Source and licence.} Sharp Endpoint Eigenfunction Estimates for the Two-Dimensional Hermite Operator. Version arXiv:2608.08423v1 (CC BY 4.0). This material is distributed under the Creative Commons Attribution 4.0 International licence, http://creativecommons.org/licenses/by/4.0/, as recorded in the arXiv OAI licence metadata for this exact version and shown on the abstract page https://arxiv.org/abs/2608.08423v1. Full rights notice: licenses/arxiv-2608.08423-CC-BY-4.0.txt.\par\smallskip
\noindent\textit{Editorial scope.} Counted unit: the original \texttt{proposition} environment \texttt{eq:package-allowed} at source lines 277--308 together with its complete proof at lines 310--377; the delivered document keeps source lines 215--1189 so that every referenced label and every citation resolves inside it. Native editable source is the paper's own concatenated TeX; the AMS wrapper is editorial formatting only. No publisher class, upstream script or unrelated artwork is redistributed.
\section{Koch--Tataru estimates}\label{sec:local}

Koch--Tataru use \(\rho^2\) for the eigenvalue \(\lambda=2N+2\). Recall that
\(u(x)=1-|x|^2/\lambda\), and set
\[
 \begin{gathered}
  \rho=\sqrt\lambda,\qquad
  \eta(x)=\frac{|x|}{\sqrt\lambda}-1,
  \qquad u_*=\lambda^{-2/3},\\
  y=\rho^{-2/3}(\rho^2-|x|^2)=\lambda^{2/3}u(x).
 \end{gathered}
\]
Thus \(u\) is the signed normalized radial coordinate relative to the turning
circle \(\{|x|=\sqrt\lambda\}\) and is positive on its interior side.
The quantity \(u_*\) is the turning-point scale. For \(\mu>0\), write
\[
 \mathcal A_\mu^{\rm in}=\{x:\tfrac12\mu\le u(x)\le2\mu\},
 \qquad
 \mathcal A_\mu^{\rm out}=\{x:\tfrac12\mu\le\eta(x)\le2\mu\}.
\]
Since
\[
 u=(1-|x|/\rho)(1+|x|/\rho),
 \qquad -y=\lambda^{2/3}(2\eta+\eta^2),
\]
each of these annuli meets only a bounded number of pieces in the
Koch--Tataru turning-point decomposition.
Choose fixed constants
\begin{equation}\label{eq:lg-constant-hierarchy}
 \begin{gathered}
  0<u_0<U_1<U_2<\frac14,\qquad 0<\eta_0<\frac14,\\
  0<\kappa_0<\kappa_1<\frac14,\qquad 0<c_1<1,
 \end{gathered}
\end{equation}
with \(u_0\) and \(\eta_0\) sufficiently small that the bounded-overlap
constants above are uniform for \(0<\mu\le u_0\) and
\(0<\mu\le\eta_0\), respectively. Choose
\(C_*\ge\max\{1,2/c_1\}\) sufficiently large, depending only on these fixed
constants, for the estimates in Sections~\ref{sec:radial}--\ref{sec:complementary}.
All implicit constants below are independent of \(N,\lambda,m,j,k\). We may
assume \(\lambda\ge\lambda_*\) for a fixed \(\lambda_*\), since the projection
norms associated with the finitely many smaller eigenvalues can be absorbed
into the implicit constants.

We first state the precise input. For
\(g\in\operatorname{Ran}\Pi_\lambda\), one has
\((\mathcal H-\rho^2)g=0\). We retain the notation
\(\ell^\infty_\rho L^p\) of Koch--Tataru for the supremum of the
\(L^p\)-norms over the pieces of their turning-point decomposition. In
dimension two, \cite[Theorem~3(a)]{KochTataru} gives, for \(2\le p\le6\),
\begin{equation}\label{eq:KT-weighted}
 \left\|
 \rho^{\,1/3-(2/3)\delta_p}
 \langle y\rangle_+^{-1/4+(5/4)\delta_p}
 \langle y\rangle_-^{1-\delta_p}g
 \right\|_{\ell^\infty_\rho L^p}
 \lesssim\norm g_2,
 \qquad
 \delta_p=\frac12-\frac1p,
\end{equation}
where \(\langle y\rangle_\pm=1+\max(\pm y,0)\).


\begin{proposition}
If \(C_*u_*\le\mu\le u_0\), then
\begin{equation}\label{eq:package-allowed}
 \norm{\one_{\mathcal A_\mu^{\rm in}}\Pi_\lambda}_{2\to10/3}
 \lesssim\lambda^{-1/10}.
\end{equation}
For \(u_*\le\mu\le u_0\),
\begin{equation}\label{eq:package-inner-L2}
 \norm{\one_{\mathcal A_\mu^{\rm in}}\Pi_\lambda}_{2\to2}
 \lesssim\mu^{1/4}.
\end{equation}
For every fixed \(A\ge1\),
\begin{equation}\label{eq:package-core}
 \norm{\one_{\{|u|\le Au_*\}}\Pi_\lambda}_{2\to10/3}
 \lesssim\lambda^{-1/10},
\end{equation}
and the fixed deep-interior region satisfies
\begin{equation}\label{eq:package-deep}
 \norm{\one_{\{u\ge u_0\}}\Pi_\lambda}_{2\to10/3}
 \lesssim\lambda^{-1/10}.
\end{equation}
If \(C_*u_*\le\mu\le\eta_0\), then
\begin{equation}\label{eq:package-outer}
 \norm{\one_{\mathcal A_\mu^{\rm out}}\Pi_\lambda}_{2\to10/3}
 \lesssim(\lambda\mu)^{-3/10},
\end{equation}
while
\begin{equation}\label{eq:package-far-outer}
 \norm{\one_{\{\eta\ge\eta_0/2\}}\Pi_\lambda}_{2\to10/3}
 \lesssim\lambda^{-1/10}.
\end{equation}
\end{proposition}

\begin{proof}
Let \(g=\Pi_\lambda f\), so \(\norm g_2\le\norm f_2\). Restricting
\eqref{eq:KT-weighted} to the three types of pieces gives
\begin{align*}
 \norm{\one_{\mathcal A_\mu^{\rm in}}g}_p
 &\lesssim
 \lambda^{-\delta_p/2}
 \mu^{1/4-(5/4)\delta_p}\norm g_2,
 &&u_*\lesssim\mu\le u_0,\\
 \norm{\one_{\{|u|\le Au_*\}}g}_p
 &\lesssim
 \lambda^{-1/6+\delta_p/3}\norm g_2,\\
 \norm{\one_{\mathcal A_\mu^{\rm out}}g}_p
 &\lesssim
 \lambda^{-5/6+\delta_p}
 \mu^{-(1-\delta_p)}\norm g_2,
 &&C_*u_*\le\mu\le\eta_0.
\end{align*}
Here the middle estimate follows because the indicated set meets only
\(O(1)\) turning-point pieces.

For \(p=10/3\), one has \(\delta_p=1/5\), and hence
\[
 -\frac{\delta_p}{2}=-\frac1{10},\qquad
 \frac14-\frac54\delta_p=0,
 \qquad
 -\frac16+\frac13\delta_p=-\frac1{10}.
\]
The first two estimates therefore prove \eqref{eq:package-allowed} and
\eqref{eq:package-core}. The first estimate with \(p=2\) proves
\eqref{eq:package-inner-L2} when \(\mu\ge C_*u_*\).
For the remaining range
\(u_*\le\mu<C_*u_*\), the middle estimate
with \(p=2\) and \(A=2C_*\) gives
\[
 \lambda^{-1/6}=u_*^{1/4}\le\mu^{1/4}.
\]
At
\(p=10/3\), the positive-\(y\) weight in \eqref{eq:KT-weighted} also has
exponent zero. Since
\(\{u\ge u_0\}\) meets only \(O(1)\) pieces, this proves
\eqref{eq:package-deep}.

The exterior estimate at \(p=10/3\) is
\[
 \norm{\one_{\mathcal A_\mu^{\rm out}}\Pi_\lambda}_{2\to10/3}
 \lesssim\lambda^{-19/30}\mu^{-4/5}
 =(\lambda\mu)^{-3/10}(\lambda^{2/3}\mu)^{-1/2}.
\]
Because \(\mu\ge C_*u_*\), it implies
\eqref{eq:package-outer}.

For the far exterior, \cite[Theorem~3(b)]{KochTataru} with \(p=\infty\)
and decay exponent one gives
\[
 \norm{\langle y\rangle_-\Pi_\lambda f}_{L^\infty(D_\rho^{\rm ext})}
 \lesssim\norm f_2,
 \qquad
 D_\rho^{\rm ext}=\{|x|>\rho+\tfrac12\rho^{-1/3}\}.
\]
For large \(\lambda\), the set \(\{\eta\ge\eta_0/2\}\) lies in
\(D_\rho^{\rm ext}\), and \(\langle y\rangle_-\gtrsim\lambda^{2/3}\)
there. Thus the corresponding \(L^2\to L^\infty\) norm is
\(O(\lambda^{-2/3})\). Interpolation with the \(L^2\) bound gives
\(\lambda^{-4/15}\le\lambda^{-1/10}\), proving
\eqref{eq:package-far-outer}. The finitely many remaining eigenvalues are
absorbed into the constant.
\end{proof}


The estimate \eqref{eq:package-outer} is summable over the exterior annuli.
Indeed, for \(g=\Pi_\lambda f\) and dyadic
\(\mu\in[C_*u_*,\eta_0]\),
\begin{align}
 \norm{\one_{\{C_*u_*\le\eta\le\eta_0\}}g}_{10/3}^{10/3}
 &\lesssim\sum_\mu
 \norm{\one_{\mathcal A_\mu^{\rm out}}g}_{10/3}^{10/3}\notag\\
 &\lesssim\lambda^{-1}\sum_\mu\mu^{-1}\norm f_2^{10/3}
 \lesssim\lambda^{-1/3}\norm f_2^{10/3}.
 \label{eq:outer-sum}
\end{align}
Here
\(\sum_\mu\mu^{-1}\lesssim(C_*u_*)^{-1}\approx\lambda^{2/3}\).
It remains to estimate the interior collar
\[
 \mathcal A_\lambda=\{x:C_*u_*\le u(x)\le u_0\}
\]
by the angular decomposition.

\section{Spectral decomposition in polar coordinates}\label{sec:polar}

Write \(\T=\R/(2\pi\mathbb Z)\). In polar coordinates
\(x=(r\cos\theta,r\sin\theta)\), with \(\theta\in\T\), set
\[
 \mathcal M_N=\{m\in\mathbb Z:|m|\le N,\ m\equiv N\pmod2\},
 \qquad
 n_m=\frac{N-|m|}{2}\in\mathbb N_0.
\]
Here \(L_n^\alpha\) denotes the generalized Laguerre polynomial of degree \(n\)
and parameter \(\alpha\), normalized as in
\cite[(1.1.37)]{Thangavelu}. For \(m\in\mathcal M_N\), let
\begin{equation}\label{eq:radial-laguerre}
 \begin{split}
  R_{N,m}(r)&=
  \left(\frac{2n_m!}{(n_m+|m|)!}\right)^{1/2}
  r^{|m|}L_{n_m}^{|m|}(r^2)e^{-r^2/2},\\
  \Phi_{N,m}(r,\theta)&=\frac1{\sqrt{2\pi}}R_{N,m}(r)e^{im\theta}.
 \end{split}
\end{equation}

\begin{lemma}[Polar Laguerre--Hermite basis]\label{lem:polar-basis}
For every \(m\in\mathcal M_N\),
\[
 \int_0^\infty|R_{N,m}(r)|^2r\dd r=1.
\]
The family
\(\{\Phi_{N,m}:m\in\mathcal M_N\}\) is an orthonormal basis of
\(\operatorname{Ran}\Pi_\lambda\), where
\(\lambda=2N+2\). Consequently, every
\(g\in\operatorname{Ran}\Pi_\lambda\) has the unique expansion
\begin{equation}\label{eq:polar-expansion}
 \begin{gathered}
  g(r,\theta)=\frac1{\sqrt{2\pi}}
  \sum_{m\in\mathcal M_N}c_mR_{N,m}(r)e^{im\theta},\\
  \sum_{m\in\mathcal M_N}|c_m|^2=\norm g_2^2.
 \end{gathered}
\end{equation}
\end{lemma}

\begin{proof}
The form of the basis is the two-dimensional case of the Hecke--Bochner
formula for Hermite projections \cite[Theorem~3.4.1]{Thangavelu}. For
completeness, we verify it directly.

Fix \(m\in\mathcal M_N\), and put \(q=|m|\) and \(k=n_m\), so that
\(2k+q=N\). In polar coordinates,
\[
 \mathcal H=-\partial_r^2-r^{-1}\partial_r-r^{-2}\partial_\theta^2+r^2.
\]
The function \(\ell(s)=L_k^q(s)\) satisfies the Laguerre equation
\[
 s\ell''(s)+(q+1-s)\ell'(s)+k\ell(s)=0.
\]
Consequently, for \(F(r)=r^qe^{-r^2/2}\ell(r^2)\), direct differentiation
gives
\[
 \left(-\partial_r^2-r^{-1}\partial_r+q^2r^{-2}+r^2\right)F
 =(4k+2q+2)F.
\]
Since \(m^2=q^2\) and \(2k+q=N\),
\[
 \mathcal H\bigl(F(r)e^{im\theta}\bigr)
 =(2N+2)F(r)e^{im\theta}.
\]
Moreover,
\[
 r^qe^{im\theta}=
 \begin{cases}
  (x_1+ix_2)^q,&m\ge0,\\
  (x_1-ix_2)^q,&m<0,
 \end{cases}
\]
so \(F(r)e^{im\theta}\) is a polynomial times \(e^{-|x|^2/2}\). It is
therefore smooth at the origin and belongs to \(\mathcal S(\R^2)\). Thus
\(\Phi_{N,m}\in\operatorname{Ran}\Pi_\lambda\).

The Laguerre orthogonality relation
\cite[(5.1.1)]{Szego} gives
\[
 \int_0^\infty s^q|L_k^q(s)|^2e^{-s}\dd s=\frac{(k+q)!}{k!}.
\]
Hence the change of variables \(s=r^2\), together with
\eqref{eq:radial-laguerre}, yields
\begin{align*}
 \int_0^\infty|R_{N,m}(r)|^2r\dd r
 &=\frac{2k!}{(k+q)!}
   \int_0^\infty r^{2q}|L_k^q(r^2)|^2e^{-r^2}r\dd r\\
 &=\frac{k!}{(k+q)!}
   \int_0^\infty s^q|L_k^q(s)|^2e^{-s}\dd s
 =1.
\end{align*}
The functions \((2\pi)^{-1/2}e^{im\theta}\) are orthonormal on \(\T\).
Therefore the functions \(\Phi_{N,m}\), \(m\in\mathcal M_N\), form an
orthonormal family in \(\operatorname{Ran}\Pi_\lambda\).

Finally, \(\mathcal M_N=\{-N,-N+2,\ldots,N\}\), so
\(\#\mathcal M_N=N+1\). By the Cartesian Hermite decomposition in the
Introduction,
\[
 \dim\operatorname{Ran}\Pi_\lambda
 =\#\{\alpha\in\mathbb N_0^2:|\alpha|=N\}=N+1.
\]
The orthonormal family therefore has the full dimension of the eigenspace and
is a basis. The unique expansion \eqref{eq:polar-expansion} and the identity
for its coefficients now follow from Parseval's identity.
\end{proof}

Set \(v_m(r)=r^{1/2}R_{N,m}(r)\). The radial Hermite equation gives
\[
 v_m''(r)+Q_m(r)v_m(r)=0,
 \qquad
 Q_m(r)=\lambda-r^2-\frac{m^2-1/4}{r^2}.
\]
For each \(m\in\mathcal M_N\), set
\begin{equation}\label{eq:nu-definition}
 \nu_m=\frac{1-\sqrt{1-4(m^2-1/4)/\lambda^2}}2,
 \qquad
 \nu_m(1-\nu_m)=\frac{m^2-1/4}{\lambda^2}.
\end{equation}
The square root in \eqref{eq:nu-definition} is real because
\(4(m^2-\tfrac14)\le4N^2-1<\lambda^2\). Moreover, \(\nu_0<0\) when
\(0\in\mathcal M_N\), whereas \(0<\nu_m<1/2\) for \(|m|\ge1\). For
\(|m|\ge1\),
\begin{equation}\label{eq:nu-m-comparable}
 \nu_m\approx\frac{m^2}{\lambda^2}.
\end{equation}
Indeed, \(1/2<1-\nu_m<1\) and
\(3m^2/4\le m^2-1/4\le m^2\).
Put \(r(u)=\sqrt{\lambda(1-u)}\). Then the exact factorization
\begin{equation}\label{eq:Q-factor-u}
 Q_m(r(u))=
 \lambda\frac{(u-\nu_m)(1-u-\nu_m)}{1-u}
\end{equation}
holds. In the collar \(u<1/4\), one has \(1-u-\nu_m>1/4\). Thus
\(Q_m(r(u))\) has the sign of \(u-\nu_m\): \(u>\nu_m\) is the
allowed region, whereas \(u<\nu_m\) is the forbidden region. For \(|m|\ge1\), \(\nu_m\)
is exactly the \(u\)-coordinate of the outer radial turning point. When
present, the value \(\nu_0<0\) is retained for uniform notation.

\section{Radial mode estimates}\label{sec:radial}

This section establishes the radial estimates needed for the angular-mode
decomposition. In the allowed region, away from the turning
point, we derive a uniform Liouville--Green representation and prove the
curvature bound for differences of the resulting phases. In the forbidden
region, we obtain rapid decay of the radial modes. 

\subsection{Liouville--Green representation away from the turning point}

\noindent\hspace*{\parindent}The Liouville--Green representation provides the
oscillatory form needed to apply Proposition~\ref{prop:weighted-sum} to the
non-glancing modes. Away from a turning point, the construction is classical;
see \cite[Chapter~6, Sections~1--5]{Olver} and \cite{OlverLG}. For completeness,
we give a detailed proof.

For \(u\in[C_*u_*,U_1]\), set
\[
 \mathcal I(u)=\{m\in\mathcal M_N:|m|\le\kappa_1\lambda\sqrt u\}.
\]
For real \(\xi\) with \(|\xi|\le\kappa_1\lambda\sqrt u\), define
\[
 Q_\xi(r)=\lambda-r^2-\frac{\xi^2-1/4}{r^2},
 \qquad
 S(u,\xi)=\int_{r(U_1)}^{r(u)}Q_\xi(s)^{1/2}\dd s.
\]
The phase is well defined. Indeed, for \(u\le v\le U_1\),
\[
 Q_\xi(r(v))
 =\lambda v-\frac{\xi^2-1/4}{\lambda(1-v)}
 \approx\lambda v>0.
\]
For \(m\in\mathcal I(u)\), put
\[
 b(u,m)=\left[\left(1-\frac{\nu_m}{u}\right)
                (1-u-\nu_m)\right]^{-1/4}.
\]

\begin{lemma}\label{lem:lg-radial}
For every \(m\in\mathcal I(U_1)\), there exist coefficients
\(d_{m,+},d_{m,-}\), depending only on \((\lambda,m)\), such that
\(|d_{m,\pm}|\le C\) and, whenever
\(u\in[C_*u_*,U_1]\) and \(m\in\mathcal I(u)\),
\begin{equation}\label{eq:lg-expansion}
 R_{N,m}(r(u))
 =\lambda^{-1/2}u^{-1/4}b(u,m)
  \sum_{\sigma=\pm}d_{m,\sigma}e^{i\sigma S(u,m)}
  +\mathcal E_m(u).
\end{equation}
Moreover,
\begin{equation}\label{eq:lg-symbol}
 |b(u,m)|\le C,
 \qquad
 |b(u,m+2)-b(u,m)|\le C(\lambda\sqrt u)^{-1},
\end{equation}
where the second estimate holds when
\(m,m+2\in\mathcal I(u)\), and
\begin{equation}\label{eq:lg-error}
 |\mathcal E_m(u)|
 \le C\lambda^{-1/2}u^{-1/4}(\lambda u^{3/2})^{-1}.
\end{equation}
\end{lemma}

\begin{proof}
Fix \(m\in\mathcal I(U_1)\), write \(p=Q_m^{1/2}\), and set
\[
 z=z(r)=\int_{r(U_1)}^r p(s)\dd s,
 \qquad
 w(z)=p(r)^{1/2}v_m(r).
\]
Direct differentiation gives
\[
 w_{zz}+w=V_mw,
 \qquad
 V_m=\frac{p''}{2p^3}-\frac{3(p')^2}{4p^4}.
\]

Let \(J=[r(U_2),r(U_1)]\). Since \(m\in\mathcal I(U_1)\), the displayed formula for
\(Q_\xi(r(v))\), with the same direct calculation for
\(U_1\le v\le U_2\), gives
\[
 |J|\approx\sqrt\lambda,
 \qquad p(r)^2\approx\lambda\quad(r\in J).
\]
Since \(v_m=r^{1/2}R_{N,m}\),
\[
 |v_m'|^2\le2r|R_{N,m}'|^2+(2r)^{-1}|R_{N,m}|^2.
\]

 Write \(R=R_{N,m}\).
The radial eigenvalue equation is
\[
 -\frac1r(rR')'+\frac{m^2}{r^2}R+r^2R=\lambda R.
\]
Multiplying by \(\overline R\,r\) and integrating by parts gives the radial energy identity
\[
 \int_0^\infty\left(|R'|^2+\frac{m^2}{r^2}|R|^2
 +r^2|R|^2\right)r\dd r
 =\lambda\int_0^\infty|R|^2r\dd r=\lambda.
\]
Here the boundary term \([rR'\overline R]_0^\infty\) vanishes. Indeed,
\eqref{eq:radial-laguerre} gives \(R=O(r^{|m|})\) and
\(R'=O(r^{|m|-1})\) near the origin when \(|m|\ge1\), while
\(R'=O(r)\) when \(m=0\); at infinity both \(R\) and \(R'\) decay
rapidly because of the Gaussian factor. The last equality uses
Lemma~\ref{lem:polar-basis}.
Together with \(\norm{v_m}_2=1\) and \(r\approx\sqrt\lambda\) on \(J\),
this yields
\[
 \int_J\bigl(|v_m'|^2+\lambda|v_m|^2\bigr)\dd r\lesssim\lambda.
\]
Hence there is \(\tilde r_m\in J\) such that
\[
 |v_m'(\tilde r_m)|\lesssim\lambda^{1/4},
 \qquad |v_m(\tilde r_m)|\lesssim\lambda^{-1/4}.
\]
Put \(\tilde z_m=z(\tilde r_m)\). On \(J\),
\(p\approx\lambda^{1/2}\) and \(|p'|\lesssim1\); therefore
\[
 |w(\tilde z_m)|+|w_z(\tilde z_m)|\lesssim1.
\]
Define
\[
 d_{m,\pm}=\frac12e^{\mp i\tilde z_m}
 \bigl(w(\tilde z_m)\mp iw_z(\tilde z_m)\bigr).
\]
Then \(|d_{m,\pm}|\lesssim1\), and Duhamel's formula is
\begin{equation}\label{eq:lg-duhamel}
 w(z)=\sum_{\sigma=\pm}d_{m,\sigma}e^{i\sigma z}
 +\int_{\tilde z_m}^z\sin(z-\zeta)V_m(\zeta)w(\zeta)\dd\zeta.
\end{equation}

Now let \(u\in[C_*u_*,U_1]\) with \(m\in\mathcal I(u)\).
Along the interval from \(\tilde r_m\) to \(r(u)\), write \(r=r(v)\).
Then \(u\le v\le U_2\), and direct differentiation gives
\[
 p(r(v))\approx\lambda^{1/2}v^{1/2},
 \qquad |p'(r(v))|\lesssim v^{-1/2},
 \qquad |p''(r(v))|\lesssim\lambda^{-1/2}v^{-3/2}.
\]
Consequently,
\[
 |V_m|\lesssim\lambda^{-2}v^{-3},
 \qquad |\dd z|=p|\dd r|\approx\lambda v^{1/2}\dd v,
\]
and hence
\[
 \int_{\tilde z_m}^{S(u,m)}|V_m(\zeta)|\dd\zeta
 \lesssim\lambda^{-1}\int_u^{U_2}v^{-5/2}\dd v
 \lesssim(\lambda u^{3/2})^{-1}.
\]
Gronwall's inequality applied to \eqref{eq:lg-duhamel} gives
\(\sup_{\tilde z_m\le z\le S(u,m)}|w(z)|\lesssim1\). Thus the integral
remainder in \eqref{eq:lg-duhamel}, evaluated at \(z=S(u,m)\), is
\(O((\lambda u^{3/2})^{-1})\).

Since \(R_{N,m}=r^{-1/2}p^{-1/2}w\), the factorization
\eqref{eq:Q-factor-u} gives
\[
 r(u)^{-1/2}p(r(u))^{-1/2}
 =\lambda^{-1/2}
  \bigl[(u-\nu_m)(1-u-\nu_m)\bigr]^{-1/4}
 =\lambda^{-1/2}u^{-1/4}b(u,m).
\]
This proves \eqref{eq:lg-expansion} and \eqref{eq:lg-error}. The reference
point \(\tilde r_m\) depends only on \((\lambda,m)\), so the same
coefficients work for every admissible \(u\).

It remains to prove \eqref{eq:lg-symbol}. If
\(k\in\mathcal I(u)\) and \(|k|\ge1\), then
\[
 0<\nu_k\le 2\frac{k^2}{\lambda^2}
 \le2\kappa_1^2u\le\frac u8.
\]
For \(k=0\), one has \(|\nu_0|\le(4\lambda^2)^{-1}\le u/8\), after
increasing the fixed lower threshold \(\lambda_*\) if necessary. Hence
\[
 1-\frac{\nu_k}{u}\approx1,
 \qquad
 1-u-\nu_k\approx1,
\]
which proves \(|b(u,k)|\lesssim1\). If
\(m,m+2\in\mathcal I(u)\), then
\[
 (\nu_{m+2}-\nu_m)(1-\nu_{m+2}-\nu_m)
 =\frac{4(m+1)}{\lambda^2}.
\]
Since \(1-\nu_{m+2}-\nu_m\approx1\) and
\(|m+1|\lesssim\lambda\sqrt u\),
\[
 |\nu_{m+2}-\nu_m|\lesssim\frac{\sqrt u}{\lambda}.
\]
On the interval joining \(\nu_m\) and \(\nu_{m+2}\), the derivative with
respect to \(\nu\) of
\(\bigl[(1-\nu/u)(1-u-\nu)\bigr]^{-1/4}\) is \(O(u^{-1})\).
The mean value theorem therefore yields
\[
 |b(u,m+2)-b(u,m)|\lesssim(\lambda\sqrt u)^{-1}.
\]
\end{proof}



\begin{lemma}[Curvature of the radial phase]\label{lem:phase-second-variation}
For \(\tau>0\), write \(u_\tau=\lambda^2\tau^2\). Let \(s,t>0\) satisfy
\[
 u_s,u_t\in[C_*u_*,U_1].
\]
Uniformly for every real \(\xi\) satisfying
\(|\xi|\le\kappa_1\lambda^2\min(s,t)\),
\[
 \left|\partial_\xi^2\bigl(S(u_t,\xi)-S(u_s,\xi)\bigr)\right|
 \approx|t-s|.
\]
If \(t\ne s\), this derivative has the sign of \(t-s\) throughout the
stated interval.
\end{lemma}

\begin{proof}
Let \(p_\xi=Q_\xi^{1/2}\). For \(v\) between \(u_s\) and \(u_t\),
the explicit formula for \(Q_\xi(r(v))\) and the bound on \(\xi\) give
\[
 p_\xi(r(v))^2\approx\lambda v,
 \qquad r(v)^2\approx\lambda,
 \qquad |\dd r|\approx\lambda^{1/2}\dd v.
\]
The common lower limit \(r(U_1)\) cancels, and the endpoints do not depend
on \(\xi\). Hence
\[
 S(u_t,\xi)-S(u_s,\xi)
 =\int_{r(u_s)}^{r(u_t)}p_\xi(r)\dd r,
 \qquad
 \partial_\xi^2p_\xi
 =-\frac1{r^2p_\xi}-\frac{\xi^2}{r^4p_\xi^3}.
\]
Since \(\xi^2/(r^2p_\xi^2)\lesssim\kappa_1^2\),
\[
 -\partial_\xi^2p_\xi\approx\frac1{r^2p_\xi}.
\]
Therefore
\[
 \left|\partial_\xi^2\bigl(S(u_t,\xi)-S(u_s,\xi)\bigr)\right|
 \approx\frac1\lambda
 \left|\int_{u_s}^{u_t}v^{-1/2}\dd v\right|
 =2|t-s|.
\]
Because \(\partial_\xi^2p_\xi<0\) and \(r(u)\) is decreasing, the derivative
has the sign of \(t-s\).
\end{proof}

\subsection{Rapid decay in the forbidden region}

If \(u\le c_1\nu_m\), then
\(\nu_m-u\ge(1-c_1)\nu_m\), so the coefficient in the rescaled equation
below is uniformly positive.

\begin{lemma}\label{lem:forbidden-radial}
Suppose that
\[
 0\le u\le c_1\nu_m,
 \qquad \nu_m\ge C_*u_*.
\]
Then, for every \(M>0\),
\begin{equation}\label{eq:forbidden-radial}
 |R_{N,m}(r(u))|
 \lesssim_M
 \lambda^{-1/2}\nu_m^{-1/4}
 \bigl(\lambda\nu_m^{3/2}\bigr)^{-M}.
\end{equation}
\end{lemma}

\begin{proof}
Put
\[
 \nu=\nu_m,\qquad \Lambda=\lambda\nu^{3/2},\qquad y=\frac u\nu,
 \qquad H(y)=(1-\nu y)^{1/4}v_m(r(\nu y)).
\]
A direct calculation from the radial equation gives
\[
 H''=\Lambda^2F_{\nu,\Lambda}(y)H,
 \qquad
 F_{\nu,\Lambda}(y)
 =\frac{(1-y)(1-\nu-\nu y)}{4(1-\nu y)^2}
 -\frac{3\nu^2}{16\Lambda^2(1-\nu y)^2}.
\]
Set \(\beta=(1+c_1)/2\). For \(y\le\beta\), write \(a=1-y\). The first
term in \(F_{\nu,\Lambda}\) satisfies
\[
 \frac{a(1-2\nu+\nu a)}{4(1-\nu+\nu a)^2}
 \ge \frac{a^2}{2(1+a)^2}\ge c>0.
\]
Indeed, after removing the factor \(a/4\), the quotient on the left is
decreasing for \(0<\nu<1/2\), and its value at \(\nu=1/2\) is
\(2a/(1+a)^2\). The second term is \(O(\Lambda^{-2})\). Since
\(\Lambda\ge C_*^{3/2}\), the choice of \(C_*\) therefore ensures that
\begin{equation}\label{eq:forbidden-positivity}
 F_{\nu,\Lambda}(y)\ge c_0>0\qquad(y\le\beta).
\end{equation}

Fix \(c_1<\gamma<\gamma'<\beta\). The localized estimate
\eqref{eq:package-inner-L2}, applied to the normalized mode
\(\Phi_{N,m}\), gives
\begin{equation}\label{eq:forbidden-mode-strip}
 \left\|\one_{\{\gamma\nu\le u\le\gamma'\nu\}}\Phi_{N,m}\right\|_2
 \lesssim \nu^{1/4}.
\end{equation}
For sufficiently small \(\nu\), this follows by covering the strip by a
bounded number of annuli \(\mathcal A_\mu^{\rm in}\); the lower restriction
\(\mu\ge u_*\) follows from \(\nu\ge C_*u_*\). For the remaining fixed range
of \(\nu\), it follows from \(\|\Phi_{N,m}\|_2=1\). Since
\[
 R_{N,m}(r(\nu y))
 =\lambda^{-1/4}(1-\nu y)^{-1/2}H(y),
 \qquad r\,\dd r=-\frac{\lambda\nu}{2}\,\dd y,
\]
\eqref{eq:forbidden-mode-strip} yields
\begin{equation}\label{eq:forbidden-strip-L2}
 \norm H_{L^2(\gamma,\gamma')}
 \lesssim(\lambda\nu)^{-1/4}.
\end{equation}

By \eqref{eq:radial-laguerre}, \(H(y),H'(y)\to0\) as \(y\to-\infty\).
Let \(W=|H|^2\). On \(( -\infty,\beta]\),
\[
 W''=2|H'|^2+2\Lambda^2F_{\nu,\Lambda}|H|^2
 \ge2c_0\Lambda^2W.
\]
Set \(\kappa=\sqrt{2c_0}\,\Lambda\). Then \(W''\ge\kappa^2W\). Since
\(W'(-\infty)=0\), one has \(W'\ge0\), and hence
\[
 \bigl((W')^2-\kappa^2W^2\bigr)'
 =2W'(W''-\kappa^2W)\ge0.
\]
The expression on the left tends to zero as \(y\to-\infty\). Thus
\(W'\ge\kappa W\), and consequently
\[
 W(y)\le e^{-\kappa(z-y)}W(z),
 \qquad y\le z\le\beta.
\]
Choose \(z\in(\gamma,\gamma')\) with
\(|H(z)|\lesssim\norm H_{L^2(\gamma,\gamma')}\). Then
\eqref{eq:forbidden-strip-L2} gives
\[
 \sup_{y\le c_1}|H(y)|
 \lesssim(\lambda\nu)^{-1/4}e^{-c\Lambda}
 \lesssim_M(\lambda\nu)^{-1/4}\Lambda^{-M}.
\]
Finally,
\[
 R_{N,m}(r(u))
 =\lambda^{-1/4}(1-u)^{-1/2}H(u/\nu).
\]
Since \(u\le c_1\nu<1/2\), this proves \eqref{eq:forbidden-radial}.
\end{proof}

\section{Weighted exponential-sum estimate}\label{sec:weighted}

This section proves the exponential-sum estimate used in the analysis of the
non-glancing modes. We first combine discrete van der Corput with Abel
summation to control exponential sums whose amplitudes have bounded variation.
A weighted $TT^*$ argument then yields Proposition~\ref{prop:weighted-sum},
which couples the radial parameter across all scales and avoids a logarithmic
loss. 

For a finite consecutive
subset \(\Gamma\) of a translate of \(2\mathbb Z\), define
\[
 \operatorname{Var}_\Gamma(a)
 =\sum_{\substack{m\in\Gamma\\m+2\in\Gamma}}|a(m+2)-a(m)|.
\]

\begin{lemma}[Weighted discrete van der Corput estimate]
\label{lem:discrete-vdc}
Let \(\Gamma\) be a finite consecutive subset of a translate of
\(2\mathbb Z\), put \(L=\#\Gamma\), and let \(0<\delta\le1\). Suppose that
\(\psi:\Gamma\to\mathbb R\) extends to a real \(C^2\) function on the convex
hull of \(\Gamma\), that \(\psi''\) has constant sign, and that
\[
 |\psi''(x)|\approx\delta.
\]
Then, for every \(a:\Gamma\to\mathbb C\),
\[
 \left|\sum_{m\in\Gamma}a(m)e^{i\psi(m)}\right|
 \lesssim\bigl(L\delta^{1/2}+\delta^{-1/2}\bigr)
 \bigl(\norm a_{\ell^\infty(\Gamma)}+\operatorname{Var}_\Gamma(a)\bigr).
\]
\end{lemma}

\begin{proof}
For every consecutive \(\Gamma_0\subset\Gamma\), reindex the parity lattice by
consecutive integers and apply the standard discrete van der Corput estimate
\cite[Chapter~2, Theorem~2.2]{GrahamKolesnik}. When
\(\#\Gamma_0\le2\), the same bound is immediate. Hence
\[
 \sup_{\Gamma_0}
 \left|\sum_{m\in\Gamma_0}e^{i\psi(m)}\right|
 \lesssim L\delta^{1/2}+\delta^{-1/2}.
\]
Abel summation on the ordered lattice proves the result.
\end{proof}

\begin{proposition}\label{prop:weighted-sum}
Let \(\lambda\ge1\), let \(I_\lambda\subset(0,\lambda^{-1})\) be an interval,
and let \(\Gamma\) be a finite consecutive subset of a translate of
\(2\mathbb Z\). Let
\(a:I_\lambda\times\Gamma\to\mathbb C\) and
\(\phi:I_\lambda\times\Gamma\to\mathbb R\) be measurable in \(t\), and set
\[
 \mathcal U(t)c(\theta)
 =\sum_{m\in\Gamma}a(t,m)e^{i(m\theta+\phi(t,m))}c_m.
\]
Assume uniformly in \(t\) that \(\operatorname{supp}a(t,\cdot)\) is contained
in a consecutive lattice interval with \(\lesssim1+\lambda^2t\) points and
\[
 \norm{a(t,\cdot)}_{\ell^\infty(\Gamma)}
 +\operatorname{Var}_\Gamma(a(t,\cdot))\lesssim1.
\]
Whenever \(t\ne s\) and the common support contains at least three points,
assume that \(\phi(t,\cdot)-\phi(s,\cdot)\) extends to a real \(C^2\) function
\(\Phi_{t,s}\) on its convex hull and
\[
 |\Phi_{t,s}''(x)|\approx|t-s|.
\]
Then, for every \(c\in\ell^2(\Gamma)\),
\begin{equation}\label{eq:weighted-sum-bound}
 \norm{t^{-1/5}\mathcal U(t)c}_{L^{10/3}(I_\lambda\times\T)}
 \lesssim\norm c_{\ell^2(\Gamma)}.
\end{equation}
\end{proposition}

\begin{proof}
Put \(p=10/3\), \(p'=10/7\), and
\(Tc(t,\theta)=t^{-1/5}\mathcal U(t)c(\theta)\). By \(TT^*\), it suffices to
prove that \(TT^*:L^{p'}(I_\lambda\times\T)\to
L^p(I_\lambda\times\T)\) is bounded.

For \(t\ne s\), the kernel of \(\mathcal U(t)\mathcal U(s)^*\) is
\[
 K_{t,s}(\omega)
 =\sum_{m\in\Gamma}a(t,m)\overline{a(s,m)}
 e^{i[m\omega+\phi(t,m)-\phi(s,m)]}.
\]
The product rule for discrete variation gives
\[
 \operatorname{Var}_\Gamma(a(t,\cdot)\overline{a(s,\cdot)})
 \le \norm{a(t,\cdot)}_\infty\operatorname{Var}_\Gamma(a(s,\cdot))
 +\norm{a(s,\cdot)}_\infty\operatorname{Var}_\Gamma(a(t,\cdot)),
\]
so the product amplitude has uniformly bounded supremum and variation. Its
support is contained in a consecutive interval with
\(\lesssim1+\lambda^2\min(t,s)\) points. If the common support has at most two
points, then \(|K_{t,s}|\lesssim1\). Otherwise, restrict the sum to its convex
hull. Since the continuous nonvanishing function \(\Phi_{t,s}''\) has constant
sign, Lemma~\ref{lem:discrete-vdc}, with \(\delta=|t-s|\), gives
\[
 |K_{t,s}(\omega)|
 \lesssim
 \bigl(1+\lambda^2\min(t,s)\bigr)|t-s|^{1/2}
 +|t-s|^{-1/2}.
\]
Because \(|t-s|\le1\) and
\[
 \lambda^2\min(t,s)|t-s|\le\lambda^2ts\le1,
\]
we obtain
\begin{equation}\label{eq:weighted-kernel}
 \norm{\mathcal U(t)\mathcal U(s)^*}_{L^1(\T)\to L^\infty(\T)}
 \lesssim|t-s|^{-1/2}.
\end{equation}
The same operator is a Fourier multiplier with uniformly bounded symbol, so
its \(L^2\to L^2\) norm is bounded. Interpolation with
\eqref{eq:weighted-kernel} yields
\begin{equation}\label{eq:osc-pp}
 \norm{\mathcal U(t)\mathcal U(s)^*}_{L^{p'}(\T)\to L^p(\T)}
 \lesssim|t-s|^{-1/5}.
\end{equation}

We also use the scalar estimate
\begin{equation}\label{eq:scalar-weighted-bound}
 \norm{\mathcal K F}_{L^p(0,\infty)}
 \lesssim\norm F_{L^{p'}(0,\infty)},
 \qquad
 \mathcal K F(t)=t^{-1/5}\int_0^\infty
 |t-s|^{-1/5}s^{-1/5}F(s)\dd s.
\end{equation}
To prove it, set
\[
 f(y)=e^{y/p'}F(e^y),\qquad
 h(x)=e^{x/p}\mathcal K F(e^x).
\]
Then \(\norm f_{p'}=\norm F_{p'}\),
\(\norm h_p=\norm{\mathcal K F}_p\), and the substitutions
\(t=e^x\), \(s=e^{x+z}\) give
\[
 h(x)=\int_\mathbb R k(z)f(x+z)\dd z,
 \qquad k(z)=e^{z/10}|1-e^z|^{-1/5}.
\]
The kernel is \(O(|z|^{-1/5})\) near zero and decays exponentially at both
ends, so \(k\in L^{5/3}(\mathbb R)\). Since
\[
 1+\frac1p=\frac1{p'}+\frac1{5/3},
\]
Young's inequality proves \eqref{eq:scalar-weighted-bound}.

Finally, Minkowski's inequality and \eqref{eq:osc-pp} give
\[
 \norm{TT^*H(t,\cdot)}_{L_\theta^p}
 \lesssim t^{-1/5}\int_{I_\lambda}|t-s|^{-1/5}s^{-1/5}
 \norm{H(s,\cdot)}_{L_\theta^{p'}}\dd s.
\]
Apply \eqref{eq:scalar-weighted-bound} after extending
\(s\mapsto\norm{H(s,\cdot)}_{L_\theta^{p'}}\) by zero outside
\(I_\lambda\). This proves the \(TT^*\) bound and hence
\eqref{eq:weighted-sum-bound}.
\end{proof}

\section{The non-glancing  modes}\label{sec:non-glancing}

We now apply the preceding weighted exponential-sum estimate to the non-glancing modes. 

For \(g\in\operatorname{Ran}\Pi_\lambda\), with coefficients as in
\eqref{eq:polar-expansion}, set
\[
 \Omega(u)=\{m\in\mathcal M_N:\nu_m\le\kappa_0^2u\}
\]
and, on \(\mathcal A_\lambda\), define
\begin{equation}\label{eq:ng-definition}
 g_{<}(r(u),\theta)
 =\frac1{\sqrt{2\pi}}\sum_{m\in\Omega(u)}
 c_mR_{N,m}(r(u))e^{im\theta}.
\end{equation}

\begin{proposition}\label{prop:below-threshold}
For every \(g\in\operatorname{Ran}\Pi_\lambda\),
\[
 \norm{g_{<}}_{L^{10/3}(\mathcal A_\lambda)}
 \lesssim\lambda^{-1/10}\norm g_2.
\]
\end{proposition}

\begin{proof}
Put \(p=10/3\). Since \(\nu_m\) is nondecreasing in \(|m|\),
\(\Omega(u)\) is a consecutive symmetric interval in \(\mathcal M_N\). For
\(m\in\Omega(u)\),
\[
 m^2=\frac14+\lambda^2\nu_m(1-\nu_m)
 \le\frac14+\kappa_0^2\lambda^2u.
\]
Hence
\begin{equation}\label{eq:below-threshold-support}
 \#\Omega(u)\lesssim1+\lambda\sqrt u,
 \qquad \Omega(u)\subset\mathcal I(u).
\end{equation}
Indeed, the second assertion follows from
\(1/4\le(\kappa_1^2-\kappa_0^2)\lambda^2u\), which is ensured by the
choice of \(C_*\). Thus Lemma~\ref{lem:lg-radial} applies to all modes in
\eqref{eq:ng-definition}.

Recall that \(u_t=\lambda^2t^2\), and put
\[
 I_\lambda=
 [\sqrt{C_*}\lambda^{-4/3},\sqrt{u_0}\lambda^{-1}]
 \subset(0,\lambda^{-1}).
\]
On this interval, \(\lambda^2t=\lambda\sqrt{u_t}\gtrsim1\). Extend \(d_{m,\sigma}\) by zero for \(m\notin\mathcal I(U_1)\), set
\(c_m^\sigma=d_{m,\sigma}c_m\), and, for \(\sigma=\pm\), define
\[
 \mathcal U_\sigma(t)h(\theta)
 =\sum_{m\in\Omega(u_t)}b(u_t,m)
 e^{i(m\theta+\sigma S(u_t,m))}h_m,
 \qquad t\in I_\lambda.
\]
Then \(\norm{c^\sigma}_2\lesssim\norm c_2\), and
Lemma~\ref{lem:lg-radial} writes the leading part of \(g_<\) as
\[
 \frac1{\sqrt{2\pi}}\lambda^{-1}t^{-1/2}
 \sum_{\sigma=\pm}\mathcal U_\sigma(t)c^\sigma(\theta).
\]

We apply Proposition~\ref{prop:weighted-sum} with
\(\Gamma=\mathcal M_N\) to each \(\mathcal U_\sigma\). Its amplitude is
\(\mathbf1_{\Omega(u_t)}(m)b(u_t,m)\). By
\eqref{eq:below-threshold-support}, its support has
\(O(1+\lambda^2t)\) points, while the two boundary jumps and
\eqref{eq:lg-symbol} give
\[
 \norm{\mathbf1_{\Omega(u_t)}b(u_t,\cdot)}_\infty
 +\operatorname{Var}_{\mathcal M_N}
   (\mathbf1_{\Omega(u_t)}b(u_t,\cdot))
 \lesssim1+\frac{\#\Omega(u_t)}{\lambda^2t}
 \lesssim1.
\]
For the formal definition of the phase on the whole lattice, extend
\(\sigma S(u_t,m)\) by zero when \(m\notin\mathcal I(u_t)\); this does not
change \(\mathcal U_\sigma(t)\), and the resulting functions are measurable.
If \(t\ne s\), then
\[
 \Omega(u_t)\cap\Omega(u_s)=\Omega(\min\{u_t,u_s\}),
\]
whose real convex hull lies in
\(\{|\xi|\le\kappa_1\lambda^2\min(t,s)\}\). On this hull the phase
difference is
\[
 \sigma\bigl(S(u_t,\xi)-S(u_s,\xi)\bigr),
\]
and Lemma~\ref{lem:phase-second-variation} gives the required second
derivative bound and constant sign. Proposition~\ref{prop:weighted-sum}
therefore applies.

The factor \((2\pi)^{-1/2}\) is harmless. Since
\(|r\dd r|=\lambda^3t\dd t\),
\[
 \begin{aligned}
 \norm{\lambda^{-1}t^{-1/2}
 \mathcal U_\sigma(t)c^\sigma}_{L^p(r\dd r\dd\theta)}=\lambda^{-1/10}
 \norm{t^{-1/5}\mathcal U_\sigma(t)c^\sigma}_{L^p(\dd t\dd\theta)}.
 \end{aligned}
\]
Thus \eqref{eq:weighted-sum-bound} bounds the displayed leading term by
\(\lambda^{-1/10}\norm c_2\).

It remains to estimate the sum of the Liouville--Green remainders. Since
\(\lambda\sqrt u\gtrsim1\) on \(\mathcal A_\lambda\), Hausdorff--Young,
H\"older, \eqref{eq:below-threshold-support}, and
\eqref{eq:lg-error},
\[
 \begin{aligned}
 \left\|
 \sum_{m\in\Omega(u)}c_m\mathcal E_m(u)e^{im\theta}
 \right\|_{L_\theta^p}\lesssim
 (\lambda\sqrt u)^{1/5}
 \lambda^{-1/2}u^{-1/4}(\lambda u^{3/2})^{-1}\norm c_2
 =\lambda^{-13/10}u^{-33/20}\norm c_2.
 \end{aligned}
\]
Using \(|r\dd r|=(\lambda/2)\dd u\) and \(u_*=\lambda^{-2/3}\), we obtain
\[
 \begin{aligned}
 \left\|
 \sum_{m\in\Omega(u)}c_m\mathcal E_m(u)e^{im\theta}
 \right\|_{L^p(\mathcal A_\lambda)}^p
\lesssim
 \lambda^{-10/3}\int_{C_*u_*}^{u_0}u^{-11/2}\dd u\,\norm c_2^p
 \lesssim\lambda^{-1/3}\norm c_2^p.
 \end{aligned}
\]
Taking the \(p\)-th root and using \(\norm c_2=\norm g_2\) completes the
proof.
\end{proof}

\section{The remaining  modes}\label{sec:complementary}


\begin{thebibliography}{99}
\bibitem[GrahamKolesnik]{GrahamKolesnik} S. W. Graham and G. Kolesnik, \emph{Van der Corput's Method of Exponential Sums}, LMS Lecture Note Series, Vol.~126, Cambridge University Press, 1991.
\bibitem[KochTataru]{KochTataru} H. Koch and D. Tataru, \emph{\(L^p\) eigenfunction bounds for the Hermite operator}, Duke Math. J. \textbf{128} (2005), no.~2, 369--392.
\bibitem[Olver]{Olver} F. W. J. Olver, \emph{Asymptotics and Special Functions}, Academic Press, New York, 1974.
\bibitem[OlverLG]{OlverLG} F. W. J. Olver, \emph{Error bounds for the Liouville--Green (or WKB) approximation}, Math. Proc. Cambridge Philos. Soc. \textbf{57} (1961), no.~4, 790--810.
\bibitem[Szego]{Szego} G. Szeg\H{o}, \emph{Orthogonal Polynomials}, fourth edition, AMS Colloquium Publications, Vol.~23, 1975.
\bibitem[Thangavelu]{Thangavelu} S. Thangavelu, \emph{Lectures on Hermite and Laguerre Expansions}, Mathematical Notes, Vol.~42, Princeton University Press, 1993.
\end{thebibliography}
\end{document}
